
% ════════════════════════════════════════
% 模型的改进与推广（08.模型改进推广.tex）
% ════════════════════════════════════════
\section{模型的改进与推广}

\subsection{模型的改进}

光伏预报误差可不再离散为低、中、高三档情景，而视为随机变量，用鲁棒优化或机会约束在目标函数中直接权衡期望购电费用与违约风险的概率上界。储能环节可区分充放电效率，并计入循环寿命损耗、自放电率与分段折旧成本，其最优充放电深度将低于本文结果，经济模型更贴近实际设备。预报误差还可结合历史气象与数值天气预报，按季节与天气类型分别估计，为各预报时刻设置差异化的储能安全裕度。

\subsection{模型的推广}

两阶段随机规划与滚动调整框架不依赖具体负荷曲线，只要求价格序列与不可控出力预测序列，可直接移植到工业园区、商业综合体、电动汽车充电站乃至虚拟电厂的能量管理。将电价序列替换为待评估的峰谷电价或尖峰电价方案，即可量化价格信号对购电曲线与储能收益的引导效果，为需求响应与储能补贴设计提供依据。本文“计划可改、改价分段”的线性化技巧也不限于电力场景，凡有预测更新且偏离计划需付出代价的问题，如电力现货市场交易、天然气购销合同与物流库存调配，均可按同样方式建模求解。
